\Askhsh{35172_SOLUTION}{1}
\begin{ylh}{1}
    \item [Ύλη από εικόνα]
\end{ylh}
\OnPar{ΛΥΣΗ}
$D_f = \mathbb{R}$, αφού $1+x^2 > 0$ για κάθε $x \in \mathbb{R}$. \\
Η $f$ είναι συνεχής στο $\mathbb{R}$ ως σύνθεση συνεχών συναρτήσεων και παραγωγίσιμη στο $\mathbb{R}$ με:
\[
f'(x) = \left( \frac{2x}{1+x^2} \right)' = \frac{(1+x^2)\cdot 2 - 2x\cdot 2x}{(1+x^2)^2} = \frac{2+2x^2-4x^2}{(1+x^2)^2} = \frac{2-2x^2}{(1+x^2)^2} = \frac{2(1-x^2)}{(1+x^2)^2}.
\]
$f'(x)=0 \Leftrightarrow 2-2x^2=0 \Leftrightarrow 2x^2=2 \Leftrightarrow x^2=1 \Leftrightarrow x=\pm 1$. \\
$f'(x)>0 \Leftrightarrow 2-2x^2>0 \Leftrightarrow 2x^2<2 \Leftrightarrow x^2<1 \Leftrightarrow -1<x<1$. \\
$f'(x)<0 \Leftrightarrow x<-1$ ή $x>1$.

Τα πρόσημα της παραγώγου της συνάρτησης $f$ φαίνονται στον ακόλουθο πίνακα:
\[
\begin{array}{|c|c|c|c|c|c|c|c|}
\hline
x & -\infty & & -1 & & 1 & & +\infty \\
\hline
f'(x) & & - & 0 & + & 0 & - & \\
\hline
f(x) & & \searrow & ~ & \nearrow & ~ & \searrow & \\
\hline
\end{array}
\]
Η συνάρτηση $f$ είναι συνεχής στο $(-\infty, 0]$ με $f'(x) < 0$ στο $(-\infty, 0)$, άρα η $f$ είναι γνησίως φθίνουσα στο $(-\infty, 0]$. \\
Η συνάρτηση $f$ είναι συνεχής στο $[0, +\infty)$ με $f'(x) > 0$ στο $(0, +\infty)$, άρα η $f$ είναι γνησίως αύξουσα στο $[0, +\infty)$. \\
Επιπλέον η συνάρτηση $f$ παρουσιάζει στο $0$ ολικό ελάχιστο το $f(0)=0$.

\bigskip

β)
\begin{align*}
f''(x) &= \left( \frac{2(1-x^2)}{(1+x^2)^2} \right)' \\
&= \frac{(2(1-x^2))'(1+x^2)^2 - 2(1-x^2)((1+x^2)^2)'}{((1+x^2)^2)^2} \\
&= \frac{-4x(1+x^2)^2 - 2(1-x^2)2(1+x^2)2x}{(1+x^2)^4} \\
&= \frac{-4x(1+x^2) - 8x(1-x^2)}{(1+x^2)^3} \\
&= \frac{-4x-4x^3-8x+8x^3}{(1+x^2)^3} = \frac{4x^3-12x}{(1+x^2)^3} = \frac{4x(x^2-3)}{(1+x^2)^3}.
\end{align*}
$f''(x)=0 \Leftrightarrow 1-x^2=0 \Leftrightarrow x=\pm 1$. \\
$f''(x)>0 \Leftrightarrow 1-x^2>0 \Leftrightarrow |x|<1 \Leftrightarrow -1<x<1$. \\
$f''(x)<0 \Leftrightarrow x<-1$ ή $x>1$.

Τα πρόσημα της δεύτερης παραγώγου της συνάρτησης $f$ φαίνονται στον ακόλουθο πίνακα:
\[
\begin{array}{|c|c|c|c|c|c|c|c|}
\hline
x & -\infty & & -1 & & 1 & & +\infty \\
\hline
f'' & & - & 0 & + & 0 & - & \\
\hline
f & & \curvearrowright & ~ & \curvearrowleft & ~ & \curvearrowright & \\
\hline
\end{array}
\]
Η συνάρτηση $f$ στρέφει τα κοίλα κάτω ή είναι κοίλη στο $(-\infty, -1]$ και στο $[1, +\infty)$ αφού είναι συνεχής σε καθένα από τα διαστήματα αυτά με $f''(x) < 0$ στο εσωτερικό τους. \\
Ενώ στο $[-1, 1]$ είναι κυρτή, αφού είναι συνεχής στο διάστημα αυτό με $f''(x) > 0$ στο εσωτερικό του. \\
Έχει σημεία καμπής τα σημεία $(-1, f(-1))$ δηλαδή το σημείο $(-1, \ln2)$ και στο σημείο $(1, f(1))$ δηλαδή το σημείο $(1, \ln2)$.